\documentclass[12pt]{article}
\usepackage[margin=1in]{geometry}
\usepackage[T1]{fontenc}
\usepackage{lmodern}
\usepackage{microtype}
\usepackage{parskip}
\setlength{\parindent}{0pt}
\setlength{\parskip}{10pt}

\begin{document}

\begin{flushright}
Wang Weijia\\
School of Economics and Finance\\
Xi'an Jiaotong University\\
Xi'an 710049, China\\
\texttt{18522704823@163.com}\\[6pt]
June 2026
\end{flushright}

\bigskip

The Editor-in-Chief\\
\textit{Journal of Data Science and Analytics}

\bigskip

Dear Editor,

We submit for your consideration the manuscript entitled
\textbf{``Dynamic Green Portfolio Optimization Using Proximal Policy
Optimization: An Analytical Framework Under Neural Network Structural
Stability.''} This paper has not been published previously, is not
under review elsewhere, and all authors have approved the submission.

\textbf{Motivation and research question.}
Deep reinforcement learning has become an active approach to dynamic
portfolio optimization, yet existing implementations restrict state
spaces to market microstructure variables and ignore firm-level
regulatory signals. China's Environmental Information Disclosure (EID)
system now scores more than 1,789 listed companies annually on a 0--100
scale; regulators levied penalties against over 11,000 A-share firms
between 2015 and 2024, making this dataset a plausible source of
policy-risk information for portfolio models. We ask whether embedding
EID scores directly into a Proximal Policy Optimization (PPO) state
space improves the structural stability of the learned policy, as
measured by the normalized Frobenius distance between networks trained
independently on non-overlapping data partitions.

\textbf{Key findings.}
Using 59,051 firm-year EID observations from 2008 to 2025, with
out-of-sample portfolio evaluation over 2022--2024, we find that adding
EID to the PPO state space reduces mean Frobenius distance by 36\%
(95\% CI: 27--46\%, $p<0.001$, Cohen's $d=1.84$). A model that
combines EID with regime-adaptive reward weighting achieves a 50\%
reduction ($d=2.01$). Critically, an ablation model that adds only
regime weighting while withholding EID produces a statistically
insignificant 3.3\% reduction ($p=0.193$), isolating the stability
gain to EID's informational content. A permutation test that shuffles
EID labels within years confirms that the effect requires the genuine
firm-level ordering in disclosure scores. On a net-return basis at a
proportional cost of 0.2\% per turnover unit, the full model achieves a
Sharpe ratio of 0.58 versus 0.29 for Markowitz minimum-variance
portfolio and 0.31 for a static ESG tilt targeting the same securities.
A mediation analysis attributes 28\% of the stability improvement to
the reward-variance channel, consistent with the information-asymmetry
mechanism proposed in the theoretical framework.

\textbf{Contribution and fit with the journal.}
The paper makes three contributions of direct interest to readers of
\textit{JDSA}. First, it introduces normalized Frobenius distance as a
data-driven diagnostic for state-variable quality in reinforcement
learning models, providing a principled criterion for deciding which
information sources merit inclusion in the state space. Second, the
controlled four-model ablation design offers a replicable template for
decomposing contributions of individual signal components in DRL
portfolio research. Third, the paper bridges the environmental
disclosure literature and the machine learning literature by showing
that a governance signal widely studied in accounting and finance
research directly improves optimization algorithm behavior. All code,
hyperparameters, and data construction notes accompany the submission;
replication is fully documented at the linked repository.

\textbf{Limitations acknowledged.}
The paper is transparent about its limitations. The full causal chain
from EID to parameter stability is not established under the
observational design; we explicitly identify the quasi-experimental
variation, specifically staggered mandatory disclosure adoption, that
would be needed for causal identification. Results are currently
calibrated to Chinese A-shares, and cross-market generalizability
remains an open question.

We believe this manuscript makes a timely and well-evidenced
contribution that is well-suited to the scope of
\textit{JDSA}. We would be grateful for the opportunity to revise the
paper in light of reviewer comments.

\bigskip

Yours sincerely,

\vspace{1.5cm}

Wang Weijia\\
School of Economics and Finance, Xi'an Jiaotong University\\
\texttt{18522704823@163.com}

\end{document}
